% Chapter 5, typeset from lectures/L05: the README is the opener and the closing review, and each
% part of Appendix A is a section, in the same order and with the same subsections. Appendix B,
% the exercises, is the last section.

\chapter{Variables and the Hardware Underneath}
\label{c5:ch}

{\large\itshape The one construct left, and what the toolchain does with all of it.\par}

\begin{chapterbox}{In this chapter}
\begin{itemize}
  \item \code{variable}: process-local storage, updated immediately where a \code{signal}
    schedules.
  \item Why the same accumulation gives two different answers depending on which you reach for.
  \item What an FPGA is physically made of: look-up tables and flip-flops, and nothing that is a
    gate.
  \item The critical path and Fmax, and the latch you infer by accident.
  \item Chapter~\ref{c4:ch}'s double-flop synchronizer extracted as a generic module
    (Exercise~\ref{c5:ex:4}), and what a generic costs in the hardware Quartus builds.
  \item \code{parity_gen} on the DE0-CV: the LED toggles whenever exactly one switch changes.
\end{itemize}
\end{chapterbox}

The course builds this chapter live, in this order. First \code{parity_gen} with a \code{signal}
accumulator, written the way a programmer would write it, then traced by hand, one loop iteration
at a time, to the wrong answer. Then \code{parity_gen} with a \code{variable}: the same loop, traced
the same way, to the right one. Then what Quartus made of it: what it actually built, how fast it
says the design will run, and an inferred-latch warning produced on purpose. And finally onto the
board, \code{parity_gen} on the DE0-CV: eight switches in, one LED out, toggling every time exactly
one switch changes.

The first step is the point of the chapter. The code looks correct, simulates without complaint,
and is wrong, and the only way to see why is to trace it rather than read it.

Nothing new is assumed beyond Chapters~\ref{c1:ch} to~\ref{c4:ch}: every construct except
\code{variable} has already appeared in a design you have written. The exercises that close the
chapter include two written exercises about what the toolchain does to your VHDL, and a code review
(Exercise~\ref{c5:ex:7}) of a design that compiles cleanly and is wrong anyway.
Appendix~\ref{app:quartus} is the toolchain reference behind the chapter's demonstrations; no
exercise requires you to have Quartus or a board.

% Sections 5.1 to 5.4, from lectures/L05/appendix/a_variables_and_hardware.md: one section per
% part, A.1 to A.4.

\section{What This Chapter Is About}
\label{c5:a1}
Every construct you need in order to describe hardware, you have already used: \code{entity} and
\code{architecture} in Chapter~\ref{c1:ch}, \code{std_logic_vector}, \code{process} and \code{case}
in Chapter~\ref{c2:ch}, the clocked process and the scheduling rule in Chapter~\ref{c3:ch},
subcomponents and generics in Chapters~\ref{c2:ch} and~\ref{c4:ch}.

Two things are left, and they are the two hardest to work out on your own:
\begin{itemize}
  \item The first is \code{variable}, the one construct the book has not needed until now:
    \begin{itemize}
      \item It looks like a signal and behaves nothing like one, and the difference is exactly the
        kind that produces a design which simulates plausibly and is quietly wrong.
      \item \Secref{c5:a2} works through the case where it bites.
    \end{itemize}
  \item The second is what any of this actually becomes:
    \begin{itemize}
      \item Since Chapter~\ref{c1:ch} you have taken it on trust that the synthesis tool builds the
        circuit you described. It does, but not out of anything you would recognize.
      \item \Secref{c5:a3} opens the box. An FPGA turns out to contain no gates at all, which is why
        Chapter~\ref{c2:ch}'s Karnaugh maps bought you less than they seemed to. Your clock has a
        speed limit, and something specific sets it. And a \code{case} with a branch missing will
        quietly hand you a memory element nobody asked for.
    \end{itemize}
\end{itemize}
Neither section answers ``is my logic right?'' You have been answering that since
Chapter~\ref{c1:ch}. They answer the two questions after it: is this what I meant, and will it fit
and run.

\section{\code{variable}: Process-Local Storage with an Immediate Update}
\label{c5:a2}
A \code{variable} is declared inside a \code{process}, between \code{process(...)} and
\code{begin}. The critical difference from a \code{signal}: variable assignment uses \code{:=} and
takes effect \textbf{immediately}, so the very next statement that reads it sees the updated value,
exactly like an assignment in C. There is no scheduling, none of the
wait-until-the-process-suspends delay that \secref{c3:a6} spelled out for signals, and no ``last
assignment wins''.

\begin{booktable}{@{}p{9.5em}LL@{}}
  & \thead{\code{signal} (\code{<=})} & \thead{\code{variable} (\code{:=})} \\
  \midrule
  Declared & Architecture's declarative part & Process's declarative part \\
  Visible from & Anywhere in the architecture & Only within its own process \\
  Update timing & Scheduled: applied when the process next suspends & Immediate: applied before
    the next statement \\
  Repeated assignment in one pass & Only the last one survives; earlier ones are discarded & Every
    one takes effect, in order \\
\end{booktable}

\subsection{Why this actually matters: an XOR-parity generator}
\label{c5:parity}
An entity computing the XOR-parity of an 8-bit input, \code{'1'} if \code{bits} has an odd number of
set bits:

\begin{vhdlcode}
entity parity_gen is
    port(bits  : in  std_logic_vector(7 downto 0);
         parity: out std_logic);
end entity;
\end{vhdlcode}

\noindent Coming from software, the first instinct is to accumulate the running XOR in a loop, using
a signal as the accumulator:

\begin{vhdlcode}
-- Don't do this: shown to demonstrate why it doesn't work.
architecture broken of parity_gen is
signal parity_s: std_logic;
begin
    parity <= parity_s;

    process(bits) is
    begin
        parity_s <= '0';
        for i in bits'range loop
            parity_s <= parity_s xor bits(i);
        end loop;
    end process;
end architecture;
\end{vhdlcode}

\noindent Trace one pass, using the scheduling rule from \secref{c3:a6}: every read of
\code{parity_s} returns whatever it held \emph{before this pass started} (call it
\code{old_parity}), and only the last assignment survives.
\begin{itemize}
  \item \code{parity_s <= '0';} schedules \code{'0'}. Not yet applied.
  \item \code{i = 7} (\code{bits'range} on \code{(7 downto 0)} runs high-to-low): reads
    \code{parity_s}, gets \code{old_parity}, because the \code{'0'} has not applied. Schedules
    \code{old_parity xor bits(7)}, discarding the \code{'0'}.
  \item \code{i = 6}: reads \code{parity_s}, again \code{old_parity}. Schedules
    \code{old_parity xor bits(6)}, discarding the previous one.
  \item \ldots and so on for every iteration \ldots
  \item \code{i = 0}, the last statement in the pass: still reads \code{old_parity}, schedules
    \code{old_parity xor bits(0)}, discarding everything before it.
  \item The process suspends and \code{parity_s} becomes \code{old_parity xor bits(0)}, the only one
    of nine scheduled assignments that survived.
\end{itemize}
The result depends only on \code{bits(0)} and on whatever \code{parity_s} happened to already hold.
Seven iterations and the initial \code{'0'} were computed and thrown away. This is not a corner
case: it is what a signal-as-accumulator \emph{always} does, because every iteration reads the same
frozen value and only the textually last assignment is kept.

There is a second thing wrong with it, which the trace does not show. The surviving assignment is
\code{parity_s <= parity_s xor bits(0);}, and \code{parity_s} is not in the sensitivity list: as
hardware that is a \textbf{combinational feedback loop}, an output fed back to its own input through
a gate with no flip-flop in the path. GHDL analyzes and simulates it happily, which is how the bug
reaches a simulation at all; Quartus reports the loop rather than building it. The lesson is the
same one either way, and the fix below removes both problems at once.

The same computation with a \code{variable}:

\begin{vhdlcode}
architecture behaviour of parity_gen is
signal parity_s: std_logic;
begin
    parity <= parity_s;

    PARITY_PROCESS: process(bits) is
        variable acc: std_logic;
    begin
        acc := '0';
        for i in bits'range loop
            acc := acc xor bits(i);
        end loop;
        parity_s <= acc;
    end process;
end architecture;
\end{vhdlcode}

\noindent Now \code{acc := '0'} applies immediately; \code{i = 7} reads \code{'0'}, XORs in
\code{bits(7)} and updates \code{acc} at once; \code{i = 6} reads that just-updated value, and so on
through every iteration. The loop ends with \code{acc} holding the true XOR of all eight bits, and
\code{parity_s <= acc;} schedules that one correct value, the only assignment to \code{parity_s} in
the pass.

This is the on-disk worked example, \file{lectures/L05/parity_gen/parity_gen.vhd}, in full:

\vhdlfile{lectures/L05/parity_gen/parity_gen.vhd}

\subsection{What that loop is not}
\label{c5:loopnot}
It is not eight steps the hardware performs one after another. This is where a C background
misleads hardest, so it is worth stating outright: a \code{for} loop whose bounds are fixed at
elaboration is \textbf{unrolled} by the synthesis tool. Those eight iterations become eight XOR
gates in a chain, all settling within one propagation delay, in the same clock cycle. There is no
counter for \code{i}, no branch, no iteration in time. The loop writes eight gates in three lines;
it does not do eight things in sequence.

That is also why \code{acc}, which reads like a running total, costs no storage at all. It never has
to survive a clock edge: each of its values is just the wire between two of those XOR gates. A
\code{variable} is the right tool precisely because it never outlives one pass.

The same goes for \code{bits'range}. The \code{'range} \textbf{attribute} gives the index range of
\code{bits}, here \code{7 downto 0}, so the loop covers every bit without writing the bounds twice.
It is resolved at elaboration, like everything else in the loop header, which is exactly why the
tool can unroll it. \code{for i in 0 to 1 loop} in Chapter~\ref{c3:ch}'s \file{led_toggle.vhd}
means the same thing: both bits get their own gates and both update in the same cycle.

\subsection{The rule of thumb}
\label{c5:rule}
\begin{itemize}
  \item Use a \code{variable} for a scratch computation consumed entirely within one pass: loop
    accumulators, temporary swaps, intermediate values handed off to a signal at the end, exactly
    like \code{acc}.
  \item Use a \code{signal} for anything another process or the outside world reads, and anything
    that must hold its value across separate passes (a counter, or a shift register's stored bits,
    both in Chapter~\ref{c6:ch}).
\end{itemize}
A variable \emph{can} retain its value between passes through the same process, behaving like
static local storage rather than a fresh local each time, so it is possible to build state with
one. But because its updates are immediate rather than scheduled, it behaves differently from a
signal used the same way, particularly for anything modelling a chain of registers stepping one
stage per clock edge. If state genuinely needs to persist and be read by other logic, reach for a
\code{signal}: the scheduled-update semantics is what makes Chapter~\ref{c6:ch}'s counters and
shift registers work.

\section{What the FPGA Actually Builds}
\label{c5:a3}

\subsection{The two things the fabric is made of}
\label{c5:fabric}
An FPGA is not a blank slate onto which gates are etched. It is a fixed grid of identical small
blocks, manufactured before anybody wrote your VHDL, and ``synthesis'' is working out how to
configure them so they behave like the circuit you described. Two matter here:
\begin{itemize}
  \item A \textbf{look-up table}, or \textbf{LUT}: a tiny memory, typically four to six inputs and
    one output. Load it with the truth table of any Boolean function of those inputs and it computes
    that function. This is the honest answer to ``what does a gate become on an FPGA'': nothing
    becomes a gate. A four-input LUT computes \code{a and b}, \code{a xor b or (c and not d)}, or
    any other function of up to four inputs at exactly the same cost, because in every case it is
    the same block holding a different table.
  \item A \textbf{flip-flop}: one bit of clocked storage, sitting next to each LUT.
\end{itemize}
On the DE0-CV's Cyclone V these come packaged together as an \textbf{ALM}, holding a fracturable
six-input LUT, usable as one six-input function or two smaller ones, plus four flip-flops. That is
the unit the Quartus fitter report counts, which is why it talks about ALMs rather than gates.

That is nearly the whole story. \code{register4} becomes four flip-flops, \code{mux_8to1} a handful
of LUTs, and Chapter~\ref{c8:ch}'s state machine LUTs feeding the flip-flops that hold the state.

It also explains something from Chapter~\ref{c2:ch} that would otherwise look like a broken
promise: the book had you minimize expressions with Karnaugh maps to save gates, then never
mentioned the saving again. On an FPGA a four-input function costs one LUT whether you minimized it
or not. Minimization still matters: it is how you understand a circuit, it is exactly right for the
discrete-logic and ASIC flows K-maps were invented for, and reducing a function below the LUT's
input width is what stops it needing a second level of LUTs. But on this board it usually buys you
nothing, and it is better to know that than to wonder.

\subsection{Why a clock has a speed limit}
\label{c5:fmax}
Signals take time to get through a LUT and along the wires between blocks. That \textbf{propagation
delay} is why a clock cannot go arbitrarily fast.

Picture the shape every design in this book has, a flip-flop, some combinational logic, another
flip-flop:
\begin{itemize}
  \item On a rising edge the first presents a new value, which has to travel through every LUT and
    wire between the two and \emph{settle} before the next edge arrives.
  \item If it has not, the second flip-flop captures a value that is still changing, and the design
    does the wrong thing reliably and mysteriously.
\end{itemize}
So the clock period has to cover three things:
\begin{itemize}
  \item The first flip-flop's \textbf{clock-to-Q} delay: the time between the edge and its output
    actually changing.
  \item The \textbf{combinational delay} through the LUTs and wires between the two.
  \item The second flip-flop's \textbf{setup time}.
\end{itemize}
The slowest such path is the \textbf{critical path}, and one over that total is the \textbf{Fmax}.

Quartus computes Fmax for you, in the timing analyzer stage of the Quartus workflow
(Appendix~\ref{app:quartus}), but only if you tell it what the clock is. That means an \code{.sdc}
constraints file with a \code{create_clock} on the \code{50 MHz} input. Without one the analyzer has
no period of yours to compare anything against, so it warns that no base clocks were constrained
and reports Fmax against a default clock instead: a number, but not one you can act on. Separately,
a design that ``fails timing'' is one where some path turned out slower than the period you did
declare.

\textbf{None of that is part of this book.} You will not write an \code{.sdc}, read a timing
report, or meet a design that fails timing, because every design here fits inside 20\,ns with room
to spare. It is worth knowing the step exists and what it needs, so that you recognize it the first
time a design of yours gets close to the limit.

\noindent Two consequences worth carrying forward:
\begin{itemize}
  \item \textbf{Deep combinational logic costs you speed, not the number of flip-flops.} Chaining a
    long expression between two registers lengthens the critical path; splitting it so part of the
    work happens on one cycle and the rest on the next shortens it. That trade, more registers for a
    faster clock, is the single most common move in high-speed digital design.
  \item \textbf{\code{50 MHz} is a 20\,ns budget.} That is the number every path in your design is
    racing against, and it is a lot of time at these speeds.
\end{itemize}

\subsection{The latch you get by accident}
\label{c5:latch}
Chapter~\ref{c3:ch} built a D latch, showed it was transparent while enabled, then set it aside in
favour of the flip-flop. Here it comes back uninvited.

\Secref{c2:a5} required a \code{when others} on every \code{case}, and gave the language's reason: a
\code{case} has to cover every value of its selector's type. The general rule behind it is that a
combinational \code{process} must assign its output on \emph{every} path through it, and that
reason is now sayable:
\begin{itemize}
  \item VHDL requires a signal to keep its old value if nothing assigns a new one.
  \item If some path through your process leaves \code{x} unassigned, you have said \code{x} must
    \emph{remember} its previous value on that path, and remembering is not something combinational
    logic can do.
  \item So the tool builds the only thing that can: a latch, exactly the one from \secref{c3:a2}.
\end{itemize}
That is an \textbf{inferred latch}, and it is almost never what anybody meant. It is a warning
rather than an error, the design will often appear to work, and the message is worth recognizing:

\begin{console}
Warning: Inferred latch(es) for signal "x"
\end{console}

\noindent The fix is never to add a latch on purpose. It is to assign the signal on every path,
which is the rule you were already following.

\section{Looking Ahead}
\label{c5:a4}
Nothing in Chapters~\ref{c6:ch} to~\ref{c8:ch} introduces a new idiom. Counters, shift registers and
timers all hold state across clock edges using \secref{c3:a6}'s clocked process and the scheduling
rule above, and Chapter~\ref{c8:ch}'s state machines add only an enumerated type, a \code{case} on
it, and one new concurrent form, the conditional signal assignment. What changes is what the state
is used for.

Two things to watch for. Chapter~\ref{c6:ch}'s \code{serial_rx8} is the first design whose
correctness you cannot check by eye on an LED, which is where the provided testbenches stop being a
formality. And \secref{c4:a8}'s generic goes to work immediately: Chapter~\ref{c7:ch}'s \code{timer}
uses one for its tick count, and \code{walking_led} instantiates that timer alongside
\code{reset_sync} and \code{button_sync}, writing almost no logic of its own.

% Section 5.5, from lectures/L05/README.md: the objectives, the questions to test yourself, and the
% next lecture.

\needspace{10\baselineskip}
\section{Review}
\label{c5:sec:review}

\subsection*{What you should be able to do now}
After this chapter, you should be able to:
\begin{itemize}
  \item Say what separates a \code{signal} from a \code{variable}, predict the different result each
    produces from the same code, and choose between them deliberately.
  \item Explain what an FPGA is physically made of, why two different-looking descriptions of the
    same circuit produce the same hardware, and why a Karnaugh map buys less here than on discrete
    logic.
  \item Explain what limits a design's clock frequency, and why adding registers can make a design
    faster rather than slower.
  \item Recognize an inferred latch from its warning, say what caused it and fix it, and say why
    ``it compiled with no errors'' promises less for an FPGA design than for a C program.
\end{itemize}

\needspace{6\baselineskip}
\subsection*{Questions to test yourself}
You should be able to explain:
\begin{enumerate}
  \item A \code{signal} is assigned twice, to different values, in one pass through a
    \code{process}. What ends up on it, and when does that update actually happen?
  \item Same question for a \code{variable}.
  \item Why does accumulating across a loop with a \code{signal} give a different, usually wrong,
    result?
  \item Nothing on an FPGA becomes a gate. What does a Boolean function become instead, and why does
    a four-input function cost the same whether or not you minimized it first?
  \item What is the critical path, and why can adding flip-flops make a design run faster?
  \item A combinational \code{case} assigns its output in three branches out of four. What does the
    synthesis tool build for the fourth, and why is that the only thing it \emph{can} build?
  \item Why is an inferred latch a warning rather than an error, and why does that make it more
    dangerous rather than less?
\end{enumerate}

\needspace{6\baselineskip}
\subsection*{Looking ahead}
Chapter~\ref{c6:ch} is about counters and shift registers:
\begin{itemize}
  \item The first designs whose state advances on its own, rather than tracking an input the way
    Chapter~\ref{c3:ch}'s and Chapter~\ref{c4:ch}'s registers do.
  \item Built from Chapter~\ref{c3:ch}'s clocked-process idiom, with this chapter's scheduling rule
    doing the real work: a shift register is a chain of signals that all update together, which is
    exactly why it shifts instead of collapsing.
\end{itemize}

% Section 5.6, from lectures/L05/appendix/b_exercises.md.

\section{Exercises}
\label{c5:sec:exercises}

\begin{aside}
\textbf{How to work these.} Each exercise is tagged with the kinds of work it asks for:
\kinds{Paper} by hand, \kinds{CircuitVerse} in the browser, \kinds{VHDL} in a module of your own.

\textbf{How to check your work.} Every exercise that asks you to write a VHDL module ships a
self-checking testbench under \file{lectures/L05/exercises} in the companion repository. Write your
module in its \file{exercises/<module>} directory, using the entity name and the \textbf{port
order} the exercise specifies, then run it with GHDL; \secref{c2:app:c} has the three commands.

No FPGA board is needed for any exercise. Your job is to get the VHDL right, and the testbench is
how you confirm it. Short answers to the parts you work on paper are in
Appendix~\ref{app:answers}; reference solutions to the VHDL are published in the companion
repository, under \file{lectures/L05/solutions}. Try each exercise yourself before looking at
either.
\end{aside}

% ------------------------------------------------------------------------------------------
\exerciseset{\code{variable} Versus \code{signal}}
\label{c5:set:1}

\exercise{Tracing the signal accumulator}{Paper}
\label{c5:ex:1}
The broken \code{parity_gen} architecture from \secref{c5:a2} is reproduced below:

\begin{vhdlcode}
architecture broken of parity_gen is
signal parity_s: std_logic;
begin
    parity <= parity_s;

    process(bits) is
    begin
        parity_s <= '0';
        for i in bits'range loop
            parity_s <= parity_s xor bits(i);
        end loop;
    end process;
end architecture;
\end{vhdlcode}

\noindent Assume:
\begin{itemize}
  \item \code{bits} has the range \code{7 downto 0}.
  \item \code{bits = "10110110"}.
  \item \code{parity_s} holds \code{'0'} before the process begins.
\end{itemize}

\begin{parts}
  \item Trace the process by hand. For the initial assignment and every loop iteration, record:
    \begin{itemize}
      \item The current index \code{i}.
      \item The value of \code{bits(i)}.
      \item The value read from \code{parity_s}.
      \item The value scheduled for \code{parity_s}.
      \item Whether that scheduled assignment survives until the process suspends.
    \end{itemize}
    Remember:
    \begin{itemize}
      \item A signal assignment does not update the signal immediately.
      \item Assignments to the same signal during one process activation replace earlier scheduled
        assignments for the same simulation time.
    \end{itemize}
  \item Determine the final value assigned to \code{parity_s} after the process suspends. Then:
    \begin{itemize}
      \item Calculate the correct XOR parity of \code{"10110110"} by hand.
      \item Compare it with the result from the broken architecture.
      \item Explain why repeatedly assigning to the signal does not create an accumulator.
    \end{itemize}
  \item Trace the variable-based \code{behaviour} architecture from the chapter's
    \file{parity_gen.vhd} using the same input. For every loop iteration, record:
    \begin{itemize}
      \item The value of \code{bits(i)}.
      \item The value of \code{acc} before the XOR operation.
      \item The value of \code{acc} immediately after the XOR operation.
    \end{itemize}
    Confirm that:
    \begin{itemize}
      \item \code{acc} retains the result of each iteration for use by the next iteration.
      \item The final value assigned to \code{parity} is the correct XOR parity.
    \end{itemize}
\end{parts}

\exercise{The smaller of two values}{VHDL}
\label{c5:ex:2}
Write an entity named \code{min_of_two}.

\noindent The entity must have:

\begin{booktable}{@{}p{3em}p{4.9em}p{14.4em}L@{}}
  \thead{Port} & \thead{Direction} & \thead{Type} & \thead{Description} \\
  \midrule
  \code{a} & in & \code{std_logic_vector(3 downto 0)} & First value. \\
  \code{b} & in & \code{std_logic_vector(3 downto 0)} & Second value. \\
  \code{m} & out & \code{std_logic_vector(3 downto 0)} & The smaller of the two. \\
\end{booktable}

\noindent Treat \code{a} and \code{b} as unsigned four-bit values.

\noindent Implement the circuit using a process:
\begin{itemize}
  \item Include \code{a} and \code{b} in the sensitivity list.
  \item Declare a process-local variable to hold the smaller value.
  \item Compare \code{a} and \code{b} as unsigned numbers.
  \item Assign the variable exactly once during each process activation.
  \item Assign the variable to \code{m} after the comparison.
\end{itemize}
Use the facilities from \code{ieee.numeric_std} when performing the unsigned comparison, following
\secref{c2:a3}.

This exercise does not require a loop. Its purpose is to use a variable as temporary process-local
storage for something \emph{other than} an accumulator: \secref{c5:a2}'s worked example builds a
value up across the iterations of a loop, which is the case a variable is usually introduced for,
and it is easy to come away thinking that is the only thing they are good for. Here the variable
simply holds an intermediate result for the few lines between computing it and using it, which is
the more common use of the two in real code.

\modulefigure{min_of_two}

\selfcheck name your entity \code{min_of_two}, with ports \code{a}, \code{b} (in) and \code{m}
(out), all \code{std_logic_vector(3 downto 0)}, declared in that order; its testbench is in
\file{exercises/min_of_two}.

\exercise{Counting set bits}{VHDL, Paper}
\label{c5:ex:3}
Write an entity named \code{ones_count8} that counts how many bits of an 8-bit input are set.

Exercise~\ref{c5:ex:1} had you trace the accumulator bug on paper and Exercise~\ref{c5:ex:2} used a
variable that was not an accumulator at all. This one is the case \secref{c5:a2} is actually about:
a value built up across the iterations of a loop. It is the only exercise in the book where you
write one.

\noindent The entity must have:

\begin{booktable}{@{}p{3em}p{4.9em}p{14.4em}L@{}}
  \thead{Port} & \thead{Direction} & \thead{Type} & \thead{Description} \\
  \midrule
  \code{bits} & in & \code{std_logic_vector(7 downto 0)} & The value to count over. \\
  \code{ones} & out & \code{natural range 0 to 8} & The number of bits of \code{bits} that are
    \code{'1'}. \\
\end{booktable}

\noindent Implement it using a single combinational process:
\begin{itemize}
  \item Put \code{bits} in the sensitivity list.
  \item Declare a process-local \code{variable} for the running count, and set it to \code{0} at the
    top of every activation. A variable keeps its value between activations, so a count that is
    never cleared keeps growing.
  \item Loop over \code{bits'range}, adding \code{1} for each bit that is set.
  \item Assign the variable to \code{ones} once, after the loop.
\end{itemize}
Then answer, in a sentence each:
\begin{itemize}
  \item Write the same thing with a \code{signal} accumulator instead, run the testbench, and record
    what it reports for the input \code{00000010}. Explain the number you get using the scheduling
    rule, not by guessing.
  \item The loop runs eight times. How many adders does this design put on the FPGA, and how many
    clock cycles does it take to produce an answer? If those two answers surprise you, reread
    \secref{c5:loopnot}.
  \item \code{ones} is a \code{natural range 0 to 8} rather than a
    \code{std_logic_vector(3 downto 0)}. Both can hold the answer. What does the range buy you, and
    what would GHDL do if your count ever left it?
\end{itemize}

\modulefigure{ones_count8}

\selfcheck name your entity \code{ones_count8}, with input \code{bits}
(\code{std_logic_vector(7 downto 0)}) and output \code{ones} (\code{natural range 0 to 8}), declared
in that order; its testbench is in \file{exercises/ones_count8} and sweeps all 256 input values.

% ------------------------------------------------------------------------------------------
\exerciseset{A Module Worth Having on Its Own}
\label{c5:set:2}

\exercise{A synchronizer module}{VHDL, Paper}
\label{c5:ex:4}
Write the double-flop synchronizer as a module of its own, named \code{sync}.

You have already built this circuit. It is the first two flip-flops of Chapter~\ref{c4:ch}'s
\code{button_sync} (Exercise~\ref{c4:ex:9}), written inline there because at that point it was the
\emph{idea} that mattered: an asynchronous input needs two flip-flops before anything else may look
at it. Here you pull those two out into a module, and the exercise is not the logic, which you know,
but everything around it.

Two things make it worth the second pass. The first is that a bare synchronizer is what most
asynchronous inputs actually need: \code{button_sync} adds an edge detector because a button press
is an event, but a serial line, a status flag from another clock domain, or a switch you only ever
read the level of needs the two flip-flops and nothing more. Chapter~\ref{c8:ch}'s serial receiver
instantiates exactly this on its \code{rx} pin. The second is that it takes \textbf{two generics},
and one of them is not a number.

\begin{booktable}{@{}p{4.2em}p{10.8em}p{3.8em}L@{}}
  \thead{Generic} & \thead{Type} & \thead{Default} & \thead{Description} \\
  \midrule
  \code{SIZE} & \code{natural range 1 to 15} & \code{1} & How many independent signals to
    synchronize, and therefore the width of both vector ports. \\
  \code{PRESET} & \code{std_logic} & \code{'0'} & The value both stages take on reset. \\
\end{booktable}

\noindent and these ports:

\begin{booktable}{@{}p{5.6em}p{4.9em}>{\raggedright\arraybackslash}p{11.2em}L@{}}
  \thead{Port} & \thead{Direction} & \thead{Type} & \thead{Description} \\
  \midrule
  \code{clock} & in & \code{std_logic} & 50\,MHz system clock. \\
  \code{reset_s2_n} & in & \code{std_logic} & Active-low, \textbf{already synchronized} reset. A
    synchronizer does not synchronize its own reset; Chapter~\ref{c4:ch}'s \code{reset_sync} does
    that. \\
  \code{async_in} & in & \code{std_logic_vector(SIZE - 1 downto 0)} & The asynchronous input or
    inputs, straight from outside the clock domain. \\
  \code{sync_out} & out & \code{std_logic_vector(SIZE - 1 downto 0)} & The same signals, safe to
    use, two rising edges later. Each bit is synchronized independently. \\
\end{booktable}

\begin{parts}
  \item Implement it:
    \begin{itemize}
      \item Two internal \code{std_logic_vector(SIZE - 1 downto 0)} signals, one per stage.
      \item One process, sensitive to \code{clock} and \code{reset_s2_n}.
      \item On reset, drive both stages to \code{PRESET}, outside the clocked branch.
      \item On a rising clock edge, shift \code{async_in} through the two stages.
      \item Drive \code{sync_out} from the second stage with a concurrent assignment.
    \end{itemize}
  \item \code{PRESET} is a \code{std_logic} generic, which is the first one in this book that is not
    a number, and it exists because the safe reset value depends on what the signal \emph{means}.
    Only the instantiating design knows that. Answer:
    \begin{itemize}
      \item \code{button_sync} synchronizes active-low buttons, so it would pass \code{'1'}. What
        would \code{'0'} claim about the buttons for the two cycles after every reset?
      \item A serial line idles high, so Chapter~\ref{c8:ch}'s receiver passes \code{'1'} too. What
        would \code{'0'} look like to a receiver watching for a start bit?
      \item Given both, why is the \emph{default} \code{'0'} rather than \code{'1'}, and what does
        that say about who is responsible for getting it right?
    \end{itemize}
  \item This is also a question about what a generic costs. \code{SIZE} is
    \code{natural range 1 to 15} and \code{PRESET} is \code{std_logic}, and both are fixed before
    synthesis runs. Using \secref{c5:a3}'s picture of what the FPGA is actually made of, say how many
    flip-flops an instance at \code{SIZE = 4} uses, how many an instance at \code{SIZE = 1} uses,
    and where in the built circuit \code{PRESET} ends up. Neither generic appears as logic anywhere;
    say what happened to them instead.
  \item Verify it with its testbench. It instantiates your module twice, once one bit wide with
    \code{PRESET = '1'} and once three bits wide with \code{PRESET = '0'}, so a design that
    hard-codes either generic fails. It checks that \code{sync_out} follows \code{async_in} after
    \textbf{exactly} two rising edges rather than one or eventually, that each bit of a vector is
    synchronized on its own, and that reset loads \code{PRESET} with no clock edge involved.
\end{parts}
\note[Tip] write the architecture for \code{SIZE = 1} in your head first, then check that every line
is already correct for a vector. The aggregate \code{(others => PRESET)} and a plain vector
assignment both work element-wise, so the generalization should cost no extra code.

\modulefigure{sync}

\selfcheck name your entity \code{sync}, with generics \code{SIZE} (\code{natural range 1 to 15},
default \code{1}) and \code{PRESET} (\code{std_logic}, default \code{'0'}), inputs \code{clock},
\code{reset_s2_n}, \code{async_in} (\code{std_logic_vector(SIZE - 1 downto 0)}) and output
\code{sync_out} (\code{std_logic_vector(SIZE - 1 downto 0)}), declared in that order; its testbench
is in \file{exercises/sync}.

\textbf{Keep this file.} Chapter~\ref{c8:ch}'s serial receiver capstone instantiates it on the
\code{rx} pin, and asks you to copy it in.

% ------------------------------------------------------------------------------------------
\exerciseset{From VHDL to Hardware}
\label{c5:set:3}
The course demonstrates the Quartus workflow on a real DE0-CV board, first with the braking
assistant back in Chapter~\ref{c1:ch}, and Appendix~\ref{app:quartus} writes it up for reference.
You don't need Quartus or a board yourself; the exercises below are about understanding what that
workflow does to the VHDL you write.

\exercise{What Quartus does to your VHDL}{Paper}
\label{c5:ex:5}
Answer each of the following, referring to Appendix~\ref{app:quartus}. One or two sentences each.

\begin{parts}
  \item A Quartus project needs a \textbf{top-level entity} selected before it will compile:
    \begin{itemize}
      \item What does the tool do differently with the top-level entity than with any other entity
        in the project?
      \item \code{min_of_two} from Exercise~\ref{c5:ex:2} and its testbench are both valid VHDL.
        Why can only one of them ever be the top-level entity of an FPGA project?
    \end{itemize}
  \item Pin assignment connects a port name to a physical pin:
    \begin{itemize}
      \item Why must pins be assigned \emph{before} compiling, rather than afterwards?
      \item What would \code{m} physically be connected to if you compiled \code{min_of_two} having
        assigned no pins at all?
    \end{itemize}
  \item Synthesis turns your architecture into gates on the chip:
    \begin{itemize}
      \item Chapter~\ref{c2:ch}'s \code{combo_logic} exercise (Exercise~\ref{c2:ex:10}) had you
        write \code{x = ab + c'd} twice: once as one expression, and once with \code{c'd} named by
        an internal signal. What should synthesis produce for each, and why is your answer the same
        for both?
      \item Which of the two would you rather read six months from now?
    \end{itemize}
  \item Compilation produces warnings as well as errors:
    \begin{itemize}
      \item Why is ``it compiled with no errors'' a much weaker statement about an FPGA design than
        ``it compiled with no errors'' about a C program?
      \item Name one class of warning you would never ignore.
    \end{itemize}
\end{parts}

\exercise{Predicting parity}{Paper, VHDL}
\label{c5:ex:6}
\code{parity_gen} from \secref{c5:a2} is this chapter's worked example, and it ships a reference
testbench at \file{lectures/L05/parity_gen/parity_gen_tb.vhd}.

\begin{parts}
  \item Before running anything, predict \code{parity} by hand for each of these inputs:
    \begin{itemize}
      \item \code{"00000000"}
      \item \code{"10110110"}
      \item \code{"11111111"}
      \item \code{"10000000"}
    \end{itemize}
  \item Run the reference testbench over the worked example and confirm your four predictions:
\begin{shell}
cd lectures/L05/parity_gen
ghdl -a --std=93 parity_gen.vhd parity_gen_tb.vhd
ghdl -e --std=93 parity_gen_tb
ghdl -r --std=93 parity_gen_tb --assert-level=error --stop-time=10ms
\end{shell}
  \item Take any one of those inputs and flip exactly one bit, leaving the other seven unchanged:
    \begin{itemize}
      \item What happens to \code{parity}?
      \item Explain why changing exactly one input bit must \emph{always} invert the XOR-parity
        result, whichever bit you pick and whatever the starting value.
      \item This is the property that makes a parity bit able to detect a single-bit transmission
        error. What kind of error does it fail to detect, and why?
    \end{itemize}
\end{parts}

\exercise{Reviewing a latched decoder}{Paper}
\label{c5:ex:7}
A colleague sends you this combinational decoder for review. It compiles, and Quartus reports no
errors.

\begin{vhdlcode}
architecture behaviour of level_decode is
begin
    process(sel) is
    begin
        case sel is
            when "00"   => leds <= "0001";
            when "01"   => leds <= "0010";
            when "10"   => leds <= "0100";
            when others => null;
        end case;
    end process;
end architecture;
\end{vhdlcode}

\begin{parts}
  \item Quartus emits \code{Warning: Inferred latch(es) for signal "leds"}. Using \secref{c5:a3},
    explain what the tool built and why. What did the VHDL ask for that combinational logic cannot
    do?
  \item The author's defence is that \code{sel} is two bits, so \code{"11"} is the only value
    \code{when others} can match, and the design never generates it. Give two separate reasons
    that argument does not save them:
    \begin{itemize}
      \item one about what \code{sel} can actually hold, from \secref{c1:a4}.
      \item one that would still apply even if \code{sel} really could only ever be \code{"00"},
        \code{"01"} or \code{"10"}.
    \end{itemize}
  \item Fix it, in two different ways:
    \begin{itemize}
      \item by changing only the \code{when others} branch.
      \item by adding a single line \emph{before} the \code{case} statement, leaving every branch
        untouched.
    \end{itemize}
    Which of the two would you rather maintain, and why? The second scales to a \code{case} with
    twenty branches; the first does not.
  \item This is a warning, not an error, and the design will often appear to work on the board. Why
    does that make it more dangerous than an error rather than less?
\end{parts}

